\documentclass[a4paper,12pt]{article}
\usepackage[utf8]{inputenc}
\usepackage[T1]{fontenc}
\usepackage{lmodern}
\usepackage{hyperref}
\title{Wsparcie J1939 w CAN Simulator GUI}
\author{Zespół CAN Simulator GUI}
\date{\today}
\begin{document}
\maketitle
\section{Wprowadzenie}
Od wersji z Faza 12 CAN Simulator GUI obsługuje protokół SAE J1939. Nowy parser identyfikatorów 29-bitowych pozwala na dekompozycję identyfikatora na PGN (Parameter Group Number), priorytet oraz adres źródłowy (Source Address).

\section{Funkcje}
\begin{itemize}
    \item \textbf{Biblioteka PGN} – wbudowana baza popularnych PGN (ponad 50 wpisów) w pliku \texttt{j1939\_pgn\_definitions.json}.
    \item \textbf{J1939 Browser} – nowa zakładka w kategorii „Protokoły” wyświetlająca ramki J1939 z nazwami PGN.
    \item \textbf{Filtrowanie} – możliwość filtrowania po numerze PGN (szesnastkowo).
    \item \textbf{Integracja ze Snifferem} – opcjonalny widok kolumn PGN, Source Address i Priority w głównym Snifferze (checkbox „J1939 View”).
\end{itemize}
\section{Testy}
Testy jednostkowe parsera znajdują się w \texttt{tests/test\_j1939.py}.
\end{document}
